The contents are hard-independent. For example, the paging is provided as a interface here and implemented in Chapter 1.

Relative sub-project:

\begin{itemize}

\item \textbf{MECOCOA} Mecocoa is a tutorial-original operating system project, following UNISYM.
\end{itemize}

Mind about \textit{mixed kernel}: The static \verb|syscall| can be made for Monolithic Kernel, and dynamic \verb|syscomm| can be made for Micro Kernel.



\chapter{Memory Management}

\textbf{Header}: memory.h, auto-inc-by stdinc.

\textbf{Module}: \verb|mcca/memoman.cpp|



\section{Paging}

\labex{Paging}
\textbf{Header}: c/system/paging.h



\chapter{Task Management}

\textbf{Header}: \verb|c/task.h|

\textbf{Module}: \verb|mcca/taskman.cpp|



\section{Inter-process Communication}



\chapter{Interrupt and Exception}

\textbf{Module}: \verb|mcca/handler.cpp|

Mecocoa has its own custom context processing mechanism, which is designed to make the most efficient use of the memory space.

For x86:

Using INT and Call-Gate.

Both Rupt and Call automatically switch to STSS.STACK0. All that is needed is to save the context and the page table.

For x64:

Only for some irrecoverable exceptions is IST used. Otherwise, the stack switching occurs in the following situations.

\begin{itemize}

\item Normally Ring0...
\item Ring1\textasciitilde{}3 Interrupt, use RSP0
\item Ring1\textasciitilde{}3 Syscall, keep user's stack. We set \verb|IA32_FMASK (0xC0000084)| IF bit, which will mask RFLAGS aka 0x200 when syscalling. We also need setup \verb|IA32_STAR (0xC0000081)|. After switching to highest stack, saving context and switching CR3, we can STI.
\end{itemize}

Interrupt: If the CR3 is not that of kernel, switch it.

Each core own its highest stack mapped over 0xFFFFFFFFC0000000。



\chapter{System Calling}

\textbf{Module}: \verb|mcca/syscall.cpp|



\chapter{File System}

\textbf{Module}: \verb|mcca/fileman*.cpp|



\chapter{IO System}

\textbf{Module}: \verb|mcca/console.cpp|

